\label{chap:83-03}

\section{Construction généalogique du continu}

\subsection{APP, TRIT et TPK : définitions directrices}

La construction part de trois objets irréductibles et causalement
ordonnés. L'objet \emph{All Possible Paths} (ci-après, \APP)
construit le support nonadique commun, son graphe de tous les chemins finis et
les deux feuillets d'évaluation somme--produit. L'unité tritique (ci-après,
\TRIT) est l'unité locale minimale de relation topologique, géométrique,
informationnelle et, une fois ordonnée, temporelle : sa signature visible possède trois
états et son relèvement conserve la retenue qu'une lecture binaire ou
résiduelle perdrait. Le \emph{Topological Phase Kernel} (ci-après, \TPK)
est la catégorie de transport qui compose ces unités en routes et conserve,
au sein du même état, feuillet, orientation, frontière, registre généalogique, mémoire, phase et
survie.

\begin{center}
\begin{tabularx}{.94\textwidth}{@{}p{.10\textwidth}Y Y@{}}
\toprule
Objet & Construction & Invariantes transportados\\
\midrule
APP & carte positive, tore fini, graphe de Cayley, chemins et feuillets
\(\Sigma,\Pi\) & marque, cellule, voisinage, somme, produit, quotient et clôture\\
TRIT & signature ternaire visible et relèvement topologique--temporel local &
résidu, retenue, feuillet, orientation, seuil et réversion\\
TPK & catégorie de routes, fibres et extensions compatibles & cylindre,
frontière, registre généalogique, mémoire, phase, holonomie et survie\\
\bottomrule
\end{tabularx}
\end{center}

La construction exacte d'APP est établie sur un segment marqué et une
carte finie ; la droite réelle apparaît ensuite comme publication limite. L'intervalle
de référence, ses marques et le chemin orienté qui les parcourt sont :
\begin{equation}
I_{10}=[0,10],\qquad
\mathcal M_{10}=I_{10}\cap\mathbb Z=\{0,1,\ldots,10\},
\qquad
\mathsf P_{10}:0\longrightarrow1\longrightarrow\cdots
\longrightarrow9\longrightarrow10,
\label{p12-83-c3-s1:eq:app-marked-segment}
\end{equation}
Les extrémités fixent le bord et les neuf marques intérieures
\(I_{10}^{\circ}\cap\mathbb Z\) forment l'alphabet positif. L'objet
mathématique exact de départ est la carte
\begin{equation}
\mathcal D_9=\{1,\ldots,9\},
\qquad \mathcal D_8=\{1,\ldots,8\},
\qquad
\lambda_9:\mathcal D_9\xrightarrow{\sim}\mathbb Z/9\mathbb Z,
\qquad \lambda_9(9)=[0]_9.
\label{p12-83-c3-s1:eq:app-positive-chart}
\end{equation}
La marque \(9\) est simultanément le dernier représentant positif et la
publication de la classe nulle. Pour tout \(n\geq1\), la racine numérique positive et
le quotient sont séparés au moyen de
\begin{equation}
\operatorname{dr}_9(n)=1+((n-1)\bmod9),
\qquad q_9(n)=\left\lfloor\frac{n-1}{9}\right\rfloor,
\qquad n=9q_9(n)+\operatorname{dr}_9(n).
\label{p12-83-c3-s1:eq:positive-digital-root-carry}
\end{equation}
Au chemin \(1\to2\to\cdots\to9\) s'ajoute
l'unique arête de clôture \(9\to1\). Le successeur et sa mémoire sont séparés :
\begin{equation}
u_+(a)=\begin{cases}a+1,&a<9,\\1,&a=9,\end{cases}
\qquad
\kappa_+(a)=\begin{cases}0,&a<9,\\1,&a=9,\end{cases}
\qquad a+1=u_+(a)+9\kappa_+(a).
\label{p12-83-c3-s1:eq:app-successor-carry}
\end{equation}
C'est pourquoi \(9\to1\) clôt la phase, augmente la mémoire et prolonge l'état
généalogique.

Le carré de la carte produit le tore fini et son graphe de mouvements :
\begin{equation}
X_9=(\mathbb Z/9\mathbb Z)^2,
\qquad
\Gamma_{\rm APP}=\operatorname{Cay}
\bigl(X_9,\{\pm e_1,\pm e_2\}\bigr),
\qquad
\operatorname{Path}(\Gamma_{\rm APP}).
\label{p12-83-c3-s1:eq:app-cayley-torus}
\end{equation}
L'identification des côtés opposés maintient le cardinal de la carte et ajoute
la loi de continuation. Le noyau de \(8^2=64\) interstices est complété par le gnomon
minimal qui publie la classe nulle dans les deux coordonnées :
\begin{equation}
\mathcal D_9^2=\mathcal D_8^2\sqcup\mathcal G_9,
\qquad
81=64+(9+8)=64+17.
\label{p12-83-c3-s1:eq:app-gnomon-public}
\end{equation}
La même loi se prolonge à des cartes de plus grande profondeur ; l'échelle change,
mais la clôture, le bord partagé et la mémoire de tour conservent leur type.

La complétion volumétrique associée s'exprime par l'identité de
strates
\begin{equation}
10^3=(9+1)^3=729+243+27+1.
\label{p12-83-c3-s1:eq:completion-unit}
\end{equation}
\input{figures/fig_app_emrh_mini}
Le volume intérieur, les faces, les arêtes et la clôture ponctuelle sont ainsi
typés comme des couches distinctes d'une même unité complétée. Sur les 81
cellules de la carte positive agissent les évaluations
\[
\Sigma,\Pi:\mathcal D_9^2\longrightarrow\mathcal D_9,
\qquad
\Sigma(i,j)=\operatorname{dr}_9(i+j),
\qquad
\Pi(i,j)=\operatorname{dr}_9(ij),
\]
Leurs résidus et quotients relevés sont définis par
\[
\begin{aligned}
R^+(i,j)&:=\Sigma(i,j),&
Q^+(i,j)&:=\frac{i+j-R^+(i,j)}9,\\
R^\times(i,j)&:=\Pi(i,j),&
Q^\times(i,j)&:=\frac{ij-R^\times(i,j)}9,
\end{aligned}
\]
et, par conséquent,
\[
\begin{aligned}
\widehat\Sigma(i,j)&:=\bigl(R^+(i,j),Q^+(i,j)\bigr),\\
\widehat\Pi(i,j)&:=\bigl(R^\times(i,j),Q^\times(i,j)\bigr),
\end{aligned}
\qquad
\widehat\Sigma,\widehat\Pi:
\mathcal D_9^2\longrightarrow\mathcal D_9\times\mathbb Z_{\geq0}.
\]
La bijection \(\lambda_9^{\times2}:\mathcal D_9^2\to X_9\) transporte
ces opérations et leurs relèvements sur le tore. APP est donc le tuple
\begin{equation}
\APP=\bigl(\mathcal D_9,\lambda_9,X_9,\Gamma_{\rm APP},
\operatorname{Path}(\Gamma_{\rm APP}),
\Sigma,\Pi,\widehat\Sigma,\widehat\Pi\bigr),
\label{p12-83-c3-s1:eq:app-full-definition-public}
\end{equation}
Le tableau constitue une représentation bidimensionnelle de ce tuple
complet. Le feuillet additif est latin ; le feuillet multiplicatif conserve
le radical ternaire de l'anneau \(\mathbb Z/9\mathbb Z\)
\begin{equation}
R_9:=\mathbb Z/9\mathbb Z,
\qquad
\mathfrak m:=([3]_9)=\{[0]_9,[3]_9,[6]_9\}\triangleleft R_9,
\qquad \mathfrak m^2=(0).
\label{p12-83-c3-s1:eq:app-radical-369}
\end{equation}
La rotation interne est l'automorphisme
\[
\rho:\mathfrak m^3\longrightarrow\mathfrak m^3,
\qquad \rho(x_1,x_2,x_3)=(x_2,x_3,x_1),
\]
dont l'orbite distinguée est
\[
([3]_9,[6]_9,[0]_9)\xrightarrow{\rho}
([6]_9,[0]_9,[3]_9)\xrightarrow{\rho}
([0]_9,[3]_9,[6]_9)\xrightarrow{\rho}
([3]_9,[6]_9,[0]_9).
\]
La réduction directe envoie \(3,6,9\) sur \(0,0,0\). La coordonnée interne
de l'idéal conserve, en revanche, le facteur extrait :
\begin{equation}
\eta_{\mathfrak m}:(\mathfrak m,+)
\xrightarrow{\;\sim\;}(\mathbb Z/3\mathbb Z,+),
\qquad
\eta_{\mathfrak m}([3a]_9):=[a]_3,
\qquad
([3]_9,[6]_9,[0]_9)
\xmapsto{\;\eta_{\mathfrak m}^{\times3}\;}
([1]_3,[2]_3,[0]_3),
\label{p12-83-c3-s1:eq:app-369-120}
\end{equation}
de sorte que le secteur \(3,6,9\) demeure entier et que sa coordonnée interne se
lit comme \(1,2,0\). Cette application appartient au radical d'APP ; la
projection temporelle de TRIT constitue le lecteur calendaire ultérieur. Les
trois classes calendaires
sont séparées comme
\[
\mathcal C_+=\{1,4,7\},\qquad
\mathcal C_-=\{2,5,8\},\qquad
\mathcal C_0=\{3,6,9\}.
\]
L'incompatibilité entre les deux feuillets devient un invariant mesurable. Sur
\(\mathscr A_3=\mathcal D_9^3\), muni de la mesure uniforme
\(\nu_3(x)=9^{-3}\), soit
\(\mathcal H_3=L^2(\mathscr A_3,\nu_3)\). Nous définissons
\[
\Sigma_3,\Pi_3:\mathscr A_3\longrightarrow\mathcal D_9,
\qquad
\Sigma_3(x)=\operatorname{dr}_9(x_1+x_2+x_3),
\qquad
\Pi_3(x)=\operatorname{dr}_9(x_1x_2x_3),
\]
et les projecteurs orthogonaux
\[
P_{\Sigma,3}:=\mathbb E_{\nu_3}[\,\cdot\mid\sigma(\Sigma_3)],
\qquad
P_{\Pi,3}:=\mathbb E_{\nu_3}[\,\cdot\mid\sigma(\Pi_3)]
\in\mathcal B(\mathcal H_3).
\]
L'énumération conjointe des \(9^3\) mots détermine alors
\[
\bigl\|[P_{\Sigma,3},P_{\Pi,3}]\bigr\|=\frac{\sqrt3}{4},
\qquad
B_c=7I+2R=9P_++5P_-,
\qquad
459-405=54.
\]
Le premier nombre réapparaît comme préfacteur électronique ; le deuxième organise
le bloc central ; le troisième enregistre le défaut global somme--produit.

\subsection{TRIT : unité locale topologique et unité de relation temporelle}

HMT fixe d'abord la signature visible
\begin{equation}
\mathrm{TRIT}_{\rm vis}=\{-1,0,+1\},
\qquad \tau\in\mathrm{TRIT}_{\rm vis}.
\label{p12-83-c3-s1:eq:trit-visible-signature}
\end{equation}
Face à l'unité binaire, le TRIT conserve trois régimes locaux. Dans la
distribution uniforme, son information et son entropie se distinguent comme
\begin{equation}
I_{\rm TRIT}=\log3\ \text{nats}=\log_2 3\ \text{bits},
\qquad S_{\rm TRIT}=k_BI_{\rm TRIT}.
\label{p12-83-c3-s1:eq:trit-information-unit}
\end{equation}
Le bit demeure l'unité binaire ; le TRIT est l'unité minimale capable de
distinguer simultanément retour, seuil et ouverture. L'égalité
\eqref{p12-83-c3-s1:eq:trit-information-unit} quantifie l'information du triplet ; le
coût \(k_B\Theta\log3\) appartient au lecteur thermique de réinitialisation ultérieur.
La mécanique dimensionnelle conserve en outre l'antécédent arithmétique de cette signature.
Pour \((i,j)\in\mathcal D_9^2\), la cellule arithmétique relevée associée est
\begin{equation}
c_{\rm MD}(i,j;\tau)=
\bigl(i,j;R^+,Q^+,R^\times,Q^\times;\tau\bigr).
\label{p12-83-c3-s1:eq:trit-complete-local-unit}
\end{equation}
Les résidus \(R^+,R^\times\) sont la publication visible des deux feuillets ;
les quotients \(Q^+,Q^\times\) conservent la mémoire exacte de la division
euclidienne ; et \(\tau\) oriente la transition. La signature visible occupe le premier
niveau ; la cellule qui la porte ajoute les antécédents arithmétiques, et l'état
enrichi complète ensuite frontière, route, phase et survie.

Il convient de séparer quatre strates. La signature visible est \(\tau\) ; le couple
résidu--retenue est son relèvement local ; \(c_{\rm MD}\) est la cellule
arithmétique d'APP munie de cette signature ; et l'état TPK ajoute extension,
frontière, registre généalogique, phase et survie. Cette hiérarchie conserve l'identité
propre de la coordonnée observée, de son relèvement, de la cellule porteuse et de
l'histoire enrichie qui la transporte.

Dans la carte ternaire équilibrée, l'entier local relevé
\(n\in\{-4,-3,\ldots,3,4\}\) se décompose de manière unique comme
\begin{equation}
\widehat n=(r,\kappa),
\qquad n=r+3\kappa,
\qquad r,\kappa\in\{-1,0,+1\}.
\label{p12-83-c3-s1:eq:trit-balanced-lift}
\end{equation}
La signature visible est \(\tau=r\) ; \(\kappa\) conserve la retenue que la projection
\(n\mapsto r\) oublierait.
La fibre au-dessus du zéro visible contient trois états distincts,
\[
0^+=(0,+1),\qquad 0^0=(0,0),\qquad 0^-=(0,-1),
\]
car \(n=-3,0,3\) possèdent le même résidu visible et des retenues distinctes. C'est
le relèvement topologique minimal de MD : une transition orientée d'APP
est munie de l'information nécessaire pour l'inverser et pour décider quelles
prolongations demeurent compatibles.

La signature possède en outre une règle de mémoire de feuillet. Si \(m_k\) est le mode
arithmétique actif, alors
\begin{equation}
m_{k+1}=
\begin{cases}
\Sigma,&\tau_k=+1,\\
\Pi,&\tau_k=-1,\\
m_k,&\tau_k=0.
\end{cases}
\label{p12-83-c3-s1:eq:trit-zero-memory}
\end{equation}
Le zéro code la rétention du feuillet précédent. C'est pourquoi le mot
\((+1,-1,0)\) publie la séquence \((\Sigma,\Pi,\Pi)\) : le troisième état
conserve le mode multiplicatif actif.

La relation temporelle s'obtient en paramétrant ces cellules de manière ordonnée.
Chaque signature porte une algèbre quadratique
\begin{equation}
\mathbb A_\tau=\mathbb R[X]/(X^2+\tau),
\qquad J_\tau=[X],
\label{p12-83-c3-s1:eq:trit-quadratic-algebra-local}
\end{equation}
de sorte que
\[
\begin{array}{c@{\qquad}c@{\qquad}l}
\toprule
\tau & J_\tau^2 & \text{relation temporelle ordonnée}\\
\midrule
+1 & -I & \text{retour elliptique, mémoire et passé},\\
0  & 0  & \text{seuil parabolique et présent},\\
-1 & +I & \text{ouverture hyperbolique bidirectionnelle et futur}.\\
\bottomrule
\end{array}
\]
Le TRIT précède le temps extérieur : il classe d'abord le type algébrique et
topologique du transport ; ensuite, un paramétrage ordonné le publie
comme relation temporelle. La double direction est typée par l'involution
exacte
\begin{equation}
\mathcal C_{\rm rev}(\tau_1,\ldots,\tau_n)
=(-\tau_n,\ldots,-\tau_1),
\qquad \mathcal C_{\rm rev}^2=I,
\label{p12-83-c3-s1:eq:trit-time-reversal}
\end{equation}
qui échange retour et ouverture et conserve le seuil. Dans le feuillet
hyperbolique, \(P_\pm=(I\pm J_{-1})/2\) séparent les deux directions propres et
la réversion échange leurs orientations. Topologie et temps sont ainsi les
deux lectures coordonnées du TRIT.

L'involution combinatoire \(\mathcal C_{\rm rev}\) et une conjugaison
opératorielle \(C\) appartiennent cependant à des domaines distincts. Dans une
réalisation dynamique séparée, munie de \(H=H^*\),
\(CJ_EC^{-1}=-J_E\) et \(CHC^{-1}=H\), nous fixons une échelle d'action interne
\(\hbar_{\rm int}>0\) et écrivons
\(h_{\rm int}:=2\pi\hbar_{\rm int}\). Alors
\begin{equation}
U(t)=\exp(-J_EHt/\hbar_{\rm int}),
\qquad CU(t)C^{-1}=U(-t).
\label{p12-83-c3-s1:eq:trit-bidirectional-propagator}
\end{equation}
L'entrelaceur qui transporte les mots TRIT vers la réalisation dynamique fixe
la correspondance entre \(\mathcal C_{\rm rev}\) et \(C\). La double
orientation locale fournit le domaine de cette construction.

\subsection{TPK : catégorie de transport et évolution de phase}

Le \emph{Topological Phase Kernel} est la catégorie de transport sur la
base observable d'APP qui compose chiffres, mots, opérateurs et relations
de prolongation dans un même état enrichi. Son architecture interne se répartit en
huit classes opératoires, différenciées par domaine, codomaine et fonction :
\begin{enumerate}[label=\textup{(\roman*)},leftmargin=2.25em]
\item support observable et état enrichi ;
\item sélecteur interne de feuillet \(M_{\rm ph}:\{1,\ldots,9\}\to\{+1,-1,0\}\) ;
\item translations cardinales et évolution visible \(F_{\rm obs}\) ;
\item lecteurs et quotients nonadique, tritique et de bloc \(1000\) ;
\item mémoire, retenue, cylindre, frontière et registre généalogique ;
\item périodicité et retour
\(C_{12}\hookrightarrow C_{108}\twoheadrightarrow C_9\),
\(\mathcal K_{27}\) et canaux \(90/120\) ;
\item prolongations \(\operatorname{Ext}_r\), cylindres, survie et
extension naturelle bidirectionnelle ;
\item foncteurs de sortie archimédienne, exceptionnelle étalonnée et dimensionnelle.
\end{enumerate}
Le retour de 27 pas est le morphisme composé
\begin{equation}
\mathcal K_{27}:=F_{\rm obs}^{27},
\qquad \mathcal K_{27}^{4}=F_{\rm obs}^{108}.
\label{p12-83-c3-s1:eq:tpk-kernel-27}
\end{equation}
Son quatrième itéré restaure la phase observable du calendrier, tandis que la
mémoire et l'état TPK complet poursuivent leur trajectoire ;
\(\mathcal K_{27}\) est le morphisme de retour interne de \(C_{108}\).
Le calendrier \(C_{108}\), la connexion nonadique, l'extension bilatérale et
les lecteurs occupent des places typées au sein de cette architecture. Soit
\(\mathsf P_{\rm APP}\) la catégorie
libre du graphe dirigé d'APP et
\(\mathsf P_{\rm APP}^{(2)}=\mathsf P_{\rm APP}\times
\mathsf P_{\rm APP}\) sa lecture simultanée somme--produit. Soit en outre
\[
\mathcal Q_{\rm obs}
=X_9^2\times\{N,E,S,O\}^2\times\Theta_{108},
\qquad \Theta_{108}=\mathbb Z/108\mathbb Z,
\]
où, pour le représentant
\(\widetilde\theta\in\{0,\ldots,107\}\),
\[
r(\theta)=1+(\widetilde\theta\bmod9),
\qquad g_0(\theta)=[\widetilde\theta]_9,
\]
\[
\tau(\theta)=
\begin{cases}
+1,&r(\theta)\in\{1,4,7\},\\
-1,&r(\theta)\in\{2,5,8\},\\
0,&r(\theta)\in\{3,6,9\},
\end{cases}
\qquad
\mu(\theta)=
\begin{cases}
\Sigma,&\tau(\theta)=+1,\\
\Pi,&\tau(\theta)=-1,\\
\varnothing,&\tau(\theta)=0.
\end{cases}
\]
L'étiquette calendaire \(\mu=\varnothing\) indique le seuil, et la mémoire
dynamique \(m_k\) de \eqref{p12-83-c3-s1:eq:trit-zero-memory} retient le dernier feuillet
actif. Une configuration
\(q=(p^\Sigma,p^\Pi,d^\Sigma,d^\Pi,\theta)\) conserve les deux positions,
les deux directions cardinales et l'indice du calendrier. L'évolution
visible \(F_{\rm obs}\) active le feuillet indiqué par TRIT, transporte la
position correspondante et fait avancer \(\theta\). Notons
\(\mathsf P_{\rm obs}\) la catégorie libre engendrée par les flèches
\(q\to F_{\rm obs}(q)\). Chaque flèche induit un couple de chemins d'APP et
définit le foncteur de base
\begin{equation}
B:\mathsf P_{\rm obs}\longrightarrow\mathsf P_{\rm APP}^{(2)}.
\label{p12-83-c3-s1:eq:tpk-bivariant-base-functor}
\end{equation}

Au-dessus de chaque \(q\) est disposée la fibre typée
\begin{equation}
\mathcal F_q=
\mathsf O_q\times\mathsf H_q\times\mathsf C_q\times\mathsf S_q
\times\mathsf B_q\times\mathsf\Lambda_q,
\label{p12-83-c3-s1:eq:tpk-fiber}
\end{equation}
Ses six facteurs conservent respectivement orientation et feuillet ; résidu et
quotient ; retenue ; signature ; cylindre et frontière ; et registre ordonné de
transport. Mémoire et survie sont des propriétés de la composition et des
relations de prolongation.

Pour une route \(r:q\to q'\), la relation
\begin{equation}
\operatorname{Ext}_r\subseteq\mathcal F_q\times\mathcal F_{q'}
\label{p12-83-c3-s1:eq:tpk-extension-relation}
\end{equation}
contient toutes les continuations admissibles et satisfait
\[
\Delta_q\subseteq\operatorname{Ext}_{\operatorname{id}_q},
\qquad
\operatorname{Ext}_{r_2}\circ\operatorname{Ext}_{r_1}
\subseteq\operatorname{Ext}_{r_2\circ r_1}.
\]
Les transports respectent les identités et la composition :
\begin{equation}
\begin{gathered}
\mathsf H(\operatorname{id}_q)=I,
\quad \mathsf C(\operatorname{id}_q)=I,
\quad \mathsf\Lambda(\operatorname{id}_q)=I,
\quad \kappa(\operatorname{id}_q)=0,\\
\mathsf H(r_2\!\circ r_1)=\mathsf H(r_2)\mathsf H(r_1),
\quad
\mathsf C(r_2\!\circ r_1)=\mathsf C(r_2)\mathsf C(r_1),
\quad
\mathsf\Lambda(r_2\!\circ r_1)=
\mathsf\Lambda(r_2)\mathsf\Lambda(r_1).
\end{gathered}
\label{p12-83-c3-s1:eq:tpk-functorial-transport}
\end{equation}
Outre ces lois fonctorielles, les coordonnées transportées obéissent à
\begin{equation}
\begin{aligned}
h'&=\mathsf H(r)(h),\\
c'&=\mathsf C(r)(c)+\kappa(r),\\
\lambda'&=\mathsf\Lambda(r)(\lambda),\\
\sigma'&=\operatorname{Sig}_{q'}(h',c',\lambda'),
\end{aligned}
\qquad
\kappa(r_2\circ r_1)
=\kappa(r_2)+\mathsf C(r_2)\kappa(r_1).
\label{p12-83-c3-s1:eq:tpk-typed-transport}
\end{equation}
La dernière identité est le cocycle de retenue qui rend associative la
composition avec mémoire.

Le TPK est alors la catégorie \(\mathsf{TPK}\) dont les objets sont
\begin{equation}
x=(q,o,h,c,\sigma,C,b,\lambda),
\qquad (o,h,c,\sigma,(C,b),\lambda)\in\mathcal F_q,
\label{p12-83-c3-s1:eq:tpk-enriched-object}
\end{equation}
et dont les morphismes \(x\to x'\) sont les triplets \((r,x,x')\) tels que
les coordonnées satisfont \eqref{p12-83-c3-s1:eq:tpk-typed-transport} et
\((f_x,f_{x'})\in\operatorname{Ext}_r\), où
\(f_x\in\mathcal F_q\) et \(f_{x'}\in\mathcal F_{q'}\). La composition est
\begin{equation}
(r_2,x',x'')\circ(r_1,x,x')
=(r_2\circ r_1,x,x''),
\label{p12-83-c3-s1:eq:tpk-morphism-composition}
\end{equation}
\Needspace{7\baselineskip}
l'identité est \((\operatorname{id}_q,x,x)\), et le foncteur
\begin{equation}
p:\mathsf{TPK}\longrightarrow\mathsf P_{\rm obs},
\qquad p(x)=q,
\qquad p(r,x,x')=r,
\label{p12-83-c3-s1:eq:tpk-base-functor}
\end{equation}
publie la configuration visible sans identifier les états d'une même fibre.
La définition fixe une catégorie au-dessus d'une base ; elle ne présuppose pas qu'il s'agisse d'une
fibration catégorique ni que toute flèche admette un relèvement cartésien.
Les projections locales sont séparées par type :
\[
\pi_{\rm TRIT}:\operatorname{Ob}(\mathsf{TPK})
\longrightarrow\mathrm{TRIT}_{\rm vis},
\qquad
\pi_{\rm TRIT}(x)=\tau(\theta),
\]
\[
\pi_{\rm loc}(x)=
\bigl(\pi_{\rm TRIT}(x),\mu(\theta),o,h,c,
\sigma,\lambda,g_0(\theta)\bigr).
\]
La première ne récupère que la signature TRIT visible. La seconde conserve
en outre l'étiquette calendaire de feuillet, l'orientation, le couple
résidu--quotient, la retenue, la signature de transport, le registre généalogique et la phase,
mais oublie encore le cylindre, son étiquette de bord et d'autres coordonnées de
\(q\). C'est pourquoi ni le TRIT visible ni son registre local
ne se substituent à l'état enrichi du TPK.

Le nom réunit trois fonctions inséparables. \emph{Topological} renvoie à
la catégorie de chemins, à l'incidence, au bord et à la prolongation par
cylindres ; \emph{Phase} renvoie à l'indice nonadique, au calendrier orienté et
à l'holonomie ; \emph{Kernel} désigne le noyau persistant de transport qui
demeure avant de contracter l'état au moyen d'un lecteur. Formellement,
l'objet est la catégorie précédente : une topologie supplémentaire sur les objets ou
les morphismes, ou son identification au noyau linéaire d'une application, exige
des données indépendantes.

En conséquence, deux routes ayant la même extrémité visible peuvent demeurer
des objets distincts par leur feuillet, quotient, retenue, frontière, cylindre, phase ou
registre. Le TPK ne choisit pas une sortie scalaire : il conserve l'antécédent commun
dont partent les publications archimédienne, exceptionnelle et dimensionnelle.

% Estas consecuencias se publican después de la construcción correlativa y
% del paso al límite de los capítulos 4 y 5. Se definen aquí para preservar
% íntegramente su procedencia, pero no se ejecutan antes de que exista el
% continuo al que se aplican.
\long\def\HMTDeferredContinuumConsequences{%
\section{Conséquences et publications du continu construit}

\subsection{Théorème d'assemblage terminal des constructions corrélatives engendrées par l'état commun}

Chaque profondeur \(k\) possède un état enrichi \(\mathcal E_k\) avec
toutes ses émissions survivantes et un unique constructeur corrélatif
\begin{equation}
\mathcal O_k:\mathcal E_k\longrightarrow\mathcal Z_k,
\qquad
\widehat x_k\longmapsto
\mathcal O\bigl(\{\mathfrak e_\alpha:
\alpha\in\operatorname{Ch}_k(\widehat x_k)\}\bigr).
\label{p12-83-c3-s1:eq:five-features-single-constructor}
\end{equation}
La valeur de \(\mathcal O_k\) conserve simultanément les coordonnées
spectrale, solénoïdale, de jauge, cohérente et déterminantale,
\begin{equation}
(\mathrm{sp},\mathrm{sol},\mathrm{gau},\mathrm{coh},\mathrm{det}),
\label{p12-83-c3-s1:eq:five-features-coordinates}
\end{equation}
comme cinq caractéristiques du même état et du même registre généalogique. Les
troncatures
\[
u^{\mathfrak e}_{\ell,k}:\mathcal E_\ell\to\mathcal E_k,
\qquad
u^{\mathcal O}_{\ell,k}:\mathcal Z_\ell\to\mathcal Z_k
\]
satisfont l'unique naturalité
\begin{equation}
u^{\mathcal O}_{\ell,k}\circ\mathcal O_\ell
=\mathcal O_k\circ u^{\mathfrak e}_{\ell,k},
\qquad \ell\ge k.
\label{p12-83-c3-s1:eq:five-features-naturality}
\end{equation}
C'est pourquoi l'action corrélative passe à la limite sous la forme
\begin{equation}
\mathcal O_\infty:
\mathcal C_{\rm cont}^{\rm disc}:=\varprojlim_k\mathcal E_k
\longrightarrow
\mathcal Z_\infty:=\varprojlim_k\mathcal Z_k.
\label{p12-83-c3-s1:eq:five-features-infinite-operator}
\end{equation}
Les cinq applications terminales sont des vues intrinsèques de cette unique valeur :
\begin{equation}
\boxed{
\mathfrak G_T^{\rm int}:=
\operatorname{View}_T\!\left(
\mathcal O_\infty(\mathcal C_{\rm cont}^{\rm disc})\right),
\quad
T\in\{\mathrm{sp},\mathrm{sol},\mathrm{gau},
\mathrm{coh},\mathrm{det}\}.}
\label{p12-83-c3-s1:eq:five-features-views}
\end{equation}
La simultanéité précède les résolutions spécialisées : chaque vue expose
la coordonnée requise et l'état complet conserve les quatre autres.

\subsection{Survie et limite projective}

Soit \(\operatorname{Surv}^{(\chi)}_N\) l'ensemble des histoires de
profondeur \(N\) qui satisfont simultanément les lois actives du canal
\(\chi\) et admettent une prolongation. Nous définissons les troncatures compatibles
\[
\tau_{N+1,N}:\operatorname{Surv}^{(\chi)}_{N+1}
\longrightarrow\operatorname{Surv}^{(\chi)}_N
\]
Celles-ci conservent le préfixe enrichi. Le continu généalogique est la limite
\begin{equation}
X_\chi=\varprojlim_N
\bigl(\operatorname{Surv}^{(\chi)}_N,\tau_{N+1,N}\bigr),
\qquad
\tau_{N+1,N}(x_{N+1})=x_N.
\label{p12-83-c3-s1:eq:inverse-survival}
\end{equation}
L'espace global des sections survivantes et l'une de ses sections
enrichies sont notés respectivement
\begin{equation}
X_\infty:=\bigsqcup_\chi X_\chi,
\qquad
\widehat\Omega_\infty\in X_\infty.
\label{p12-83-c3-s1:eq:global-survival-space}
\end{equation}
Le lecteur archimédien global est
\begin{equation}
P_{\rm arq}:=\operatorname{ev}_{\mathbb R}:
X_\infty\longrightarrow\mathbb R;
\label{p12-83-c3-s1:eq:archimedean-reader}
\end{equation}
sa restriction à chaque canal évalue l'intersection réduite à un point des cylindres
compatibles de la section correspondante.
Lorsque le lecteur exceptionnel intervient, \(X_\infty^{\rm cal}\subseteq
X_\infty\) désigne le sous-espace dans lequel le type d'étalonnage
de \eqref{p12-83-c3-s1:eq:exceptional-calibration} a été fixé.
La survie détermine la section par compatibilité interne : chaque
histoire doit pouvoir être prolongée dans le système complet. Le lecteur numérique
intervient ensuite pour certifier et publier l'intersection du
descendant déjà compatible.

\subsection{Connexion nonadique conjointe, holonomie et mémoire verticale}

À la profondeur \(K\), l'holonomie du retour nonadique est initialement la
correspondance typée
\begin{equation}
\Gamma_{9,K}=\operatorname{Hol}_{C_9}(\nabla_{\rm TPK}):
\widehat{\mathcal X}_K\rightrightarrows
\widehat{\mathcal X}_{K+9}.
\label{p12-83-c3-s1:eq:gamma9}
\end{equation}
Ses itérés sont des compositions relationnelles. Ce n'est que lorsque chaque état possède
une unique prolongation compatible qu'elle se restreint à une application, et ce n'est que lorsque cette
application est inversible qu'elle constitue une monodromie. Cette distinction conserve toutes
les extensions survivantes avant de sélectionner une section.
Après neuf phases, la phase locale revient, tandis que l'état complet avance :
\[
g(K+9)=g(K),\qquad q_{{\rm mem},K+9}=q_{{\rm mem},K}+1,
\qquad \widehat x_{K+9}\ne\widehat x_K\ \text{en général}.
\]
Bloc, cylindre, préfixe, retenue et mémoire continuent d'évoluer. Telle
est la différence entre périodicité visible et retour avec biographie.

L'état transporté au niveau (K) peut être présenté comme
\begin{equation}
\widehat x_K=
(B_K,C_K,p_K,c_K,\Sigma_K,W_K,g(K),q_{{\rm mem},K},s_K),
\label{p12-83-c3-s1:eq:enriched-state-k}
\end{equation}
où chaque coordonnée possède une loi de transport propre. Le retour de phase
est représenté sur le catalogue
\[
\mathcal U_6^{\rm cat}\cong\mathfrak S_+\times\mathfrak S_\times,
\qquad |\mathfrak S_+|=18,
\qquad |\mathfrak S_\times|=26.
\]
Les moyennes conditionnelles
\[
(E_+f)(s,e)=\frac1{18}\sum_{s'\in\mathfrak S_+}f(s',e),
\qquad
(E_\times f)(s,e)=\frac1{26}\sum_{e'\in\mathfrak S_\times}f(s,e')
\]
commutent. Le mot de phase
\(\tau=(+1,-1,0,+1,-1,0,+1,-1,0)\) applique
\(E_+,E_\times,I\) dans cet ordre, et un tour se clôt au moyen de
\begin{equation}
\mathcal K_{\rm ph}^{\,9}
=(E_+E_\times)\otimes I_{\ell^2(C_9)},
\qquad
\bigl\langle\delta_u,E_+E_\times\delta_v\bigr\rangle
=\frac1{468}\quad(u,v\in\mathcal U_6^{\rm cat}).
\label{p12-83-c3-s1:eq:joint-nine-holonomy}
\end{equation}
Les feuillets somme et produit n'agissent pas comme des filtres indépendants : ce sont deux
composantes d'une unique holonomie qui ne se clôt qu'après avoir parcouru le
cycle complet.

La colligation compression--expression sépare ce qui revient de ce qui demeure
comme mémoire. Si (J) plonge l'espace des histoires dans l'espace de
phase, si (T) est la compression survivante et \(\eta\) le canal de défaut,
alors
\begin{equation}
\boxed{
U_{\Gamma_9}J=JT+\eta,
\qquad J^*\eta=0,
\qquad I-T^*T=\eta^*\eta.}
\label{p12-83-c3-s1:eq:compression-expression}
\end{equation}
La première égalité décompose le transport en retour comprimé et
expression de mémoire ; la deuxième prouve leur orthogonalité ; la troisième mesure
exactement l'information qui ne tient pas dans l'ombre périodique. L'affirmation
centrale est ainsi formalisée : c'est la phase qui revient, non l'état complet.

La mémoire verticale s'organise au moyen de l'extension cyclique orientée
\begin{equation}
0\longrightarrow C_{12}\xrightarrow{\,\iota\,}C_{108}
\xrightarrow{\,p\,}C_9\longrightarrow0,
\qquad
c(r,r')=\left[\left\lfloor\frac{r+r'}9\right\rfloor\right]_{12}.
\label{p12-83-c3-s1:eq:c108-extension}
\end{equation}
Ici \(\iota([b]_{12})=[9b]_{108}\) et
\(p([\theta]_{108})=[\theta]_9\). L'extension n'est pas scindée :
\(C_{108}\) est d'exposant 108, tandis que
\(C_{12}\times C_9\) est d'exposant 36.
Le cocycle enregistre la retenue de chaque tour de la base \(C_9\) dans la
mémoire dodécaphasique \(C_{12}\). \(C_{108}\) est donc un calendrier
orienté de phase et de mémoire. Ce n'est qu'après la carte dimensionnelle que cette
extension se réalise comme une horloge quantique finie de 108 pas.

Le relèvement local de MD qui porte la signature TRIT est l'unité
topologique--temporelle de ce cristal généalogique.
Sur
\(\mathcal H_{\rm clk}=\mathbb C^{108}\), de base
\(\{|j\rangle:j\in\mathbb Z/108\mathbb Z\}\), soient
\[
\zeta=e^{2\pi\ii/108},\qquad
|\widetilde k\rangle=\frac1{\sqrt{108}}
\sum_{j=0}^{107}\zeta^{jk}|j\rangle,
\qquad U_{\rm clk}|j\rangle=|j+1\rangle.
\]
\Needspace{12\baselineskip}
Avec \(\omega_0=2\pi/(108t_0)\), le hamiltonien
\begin{equation}
H_{\rm clk}=\hbar_{\rm int}\omega_0
\sum_{k=0}^{107}k\,|\widetilde k\rangle\langle\widetilde k|
\label{p12-83-c3-s1:eq:trit-time-crystal-hamiltonian}
\end{equation}
engendre exactement
\begin{equation}
U_{\rm step}=\exp\!\left(
-\frac{\ii H_{\rm clk}t_0}{\hbar_{\rm int}}\right)=U_{\rm clk},
\qquad U_{\rm step}^{108}=I,
\label{p12-83-c3-s1:eq:trit-time-crystal-period}
\end{equation}
Si \(Z_{\rm clk}|j\rangle=\zeta^j|j\rangle\), alors
\begin{equation}
\langle j_0|U_{\rm step}^{-n}Z_{\rm clk}U_{\rm step}^{n}|j_0\rangle
=\zeta^{j_0+n},
\label{p12-83-c3-s1:eq:trit-time-crystal-observable}
\end{equation}
dont l'orbite possède une période primitive de 108. Dans le vocabulaire HMT, il s'agit du
cristal temporel généalogique interne : le calendrier observable se clôt,
\(C_{12}\) enregistre le tour et l'état TPK enrichi ne revient pas nécessairement.
La réalisation ultérieure comme phase physique macroscopique exige,
en outre, un protocole de Floquet, une réponse sous-harmonique primitive, la robustesse,
la persistance et un mécanisme de protection ; cette carte ne modifie pas le théorème
fini de périodicité unitaire.

\subsection{Publication archimédienne, nombre généalogique et mesure}

La publication archimédienne utilise des cylindres ternaires semi-ouverts
\begin{equation}
I_{N,L}=\left[\frac{N}{3^L},\frac{N+1}{3^L}\right),
\qquad \operatorname{diam}I_{N,L}=3^{-L}\longrightarrow0.
\label{p12-83-c3-s1:eq:ternary-cylinders}
\end{equation}
Une section compatible détermine des cylindres emboîtés et, par complétude, un
unique point réel. Le nombre HMT est la section elle-même,
\begin{equation}
x=(x_N)_N\in X_\chi
=\varprojlim_N\operatorname{Surv}^{(\chi)}_N,
\label{p12-83-c3-s1:eq:hmt-number-as-section}
\end{equation}
non le tuple formé par elle et ses évaluations. La valeur apparaît lorsqu'un
lecteur oublie une partie de la généalogie.
La précision est donc profondeur : mesurer signifie descendre à un
cylindre plus fin tandis que le préfixe causal demeure dans l'antécédent.
La droite réelle est reconstruite comme quotient archimédien d'un espace
d'histoires beaucoup plus riche.

\subsection{Réalisation dimensionnelle de l'état HMT : application d'écran, grandeur observable et carte spectrale}

La mécanique dimensionnelle commence là où l'état généalogique se couple à une
frontière de réalisation. Un corps \(K_n\) étant fixé, soit \(M_n(x)\) le module
engendré par les canaux de prolongation d'un état \(x\), soit
\(\mathcal R_n(x)\) le sous-module de ses relations et soit \(\beta_n(x)\) la
donnée de bord conservée. L'écran dimensionnel est
\begin{equation}
\Sigma_n(x)=\bigl(K_n,M_n(x),\beta_n(x),\mathcal R_n(x)\bigr),
\qquad
d_{\Sigma_n}(x)=\dim_{K_n}\frac{M_n(x)}{\mathcal R_n(x)}.
\label{p12-83-c3-s1:eq:md-screen}
\end{equation}
Si \(M_n(x)\simeq K_n^{g_n}\) et si les colonnes de \(R_n(x)\) engendrent les
relations, alors
\begin{equation}
\boxed{d_{\Sigma_n}(x)=g_n-\operatorname{rank}_{K_n}R_n(x).}
\label{p12-83-c3-s1:eq:md-rank}
\end{equation}
La dimension n'est pas adjointe comme étiquette : c'est le nombre de canaux
indépendants qui survivent après l'imposition des relations de frontière.

Une grandeur physique est une section d'une droite dimensionnelle \(L_D\) ; une
unité est une carte \(\chi_u:L_D\to\mathbb R\) qui publie sa coordonnée.
C'est pourquoi
\begin{equation}
\text{état TPK}\longrightarrow\Sigma_n(x)
\longrightarrow s_D\in\Gamma(L_D)
\xrightarrow{\;\chi_u\;}[s_D]_u\in\mathbb R
\label{p12-83-c3-s1:eq:md-measurement-chain}
\end{equation}
sépare exactement structure, type dimensionnel et nombre mesuré. Mesurer
consiste à coupler une section à un appareil réversible, à imposer la
compatibilité de bord et à appliquer le lecteur ; changer d'unités change
\(\chi_u\), non la section ni sa généalogie.

\subsection{Double projection mathématique et réalisation dimensionnelle}

La branche exceptionnelle conserve des données de jauge dont la publication réelle n'a pas
besoin. Soit \(\mathscr C_{\rm exc}\) l'espace des étalonnages et soit
\begin{equation}
\mathfrak E:
X_\infty^{\rm cal}\times\mathscr C_{\rm exc}
\longrightarrow\mathbf{Exc}
\label{p12-83-c3-s1:eq:exceptional-calibrated-reader}
\end{equation}
le lecteur d'incidence. Nous fixons une fois
\begin{equation}
\begin{aligned}
\chi_{\rm exc}=\bigl(&\text{ordre projectif},\ \text{calibre de Paley},
 \ \text{dodécade},\\
 &\text{alignement d’incidence},\ \text{chambre réticulaire}\bigr)
 \in\mathscr C_{\rm exc}
\end{aligned}
\label{p12-83-c3-s1:eq:exceptional-calibration}
\end{equation}
et abrégeons
\(\mathfrak E_{\chi_{\rm exc}}(x):=\mathfrak E(x,\chi_{\rm exc})\).
La jauge initiale est transportée par l'histoire ; elle n'est pas choisie de nouveau à
chaque profondeur.

La même limite étalonnée admet alors deux lecteurs de codomaines
distincts :
\begin{equation}
\mathbb R
\xleftarrow{\;\operatorname{ev}_{\mathbb R}\;}
X_\infty^{\rm cal}
\xrightarrow{\;\mathfrak E_{\chi_{\rm exc}}\;}
\mathbf{Exc}.
\label{p12-83-c3-s1:eq:double-projection}
\end{equation}
La première branche contracte les cylindres jusqu'à une valeur continue. La seconde
conserve incidences, orientations, jauge et frontière. Toutes deux descendent en
parallèle du même état source. Cette bifurcation explique conjointement la
continuité métrique et l'émergence de symétries exceptionnelles. La seconde
branche est une reconstruction étalonnée : elle est exacte une fois
\(\chi_{\rm exc}\) fixé, mais l'état nu ne sélectionne pas rétrospectivement
les cinq données qui composent cette jauge.

La réalisation dimensionnelle ne sérialise pas ces deux branches. C'est un troisième lecteur
du même antécédent :
\begin{equation}
P_{\rm dim}:X_\infty\longrightarrow
\mathcal R_{\rm MD},
\qquad
\mathcal R_{\rm MD}=
\{\text{action, masse, champ, géométrie, information}\}.
\label{p12-83-c3-s1:eq:dimensional-publication}
\end{equation}
De cette manière, un nombre, une symétrie et une grandeur physique peuvent partager
une généalogie sans se confondre ni se réduire les uns aux autres.

La thèse du continu est ainsi formulée avec précision : la réalité
archimédienne est l'ombre de rayon nul d'histoires discrètes infiniment
raffinables ; ce que l'ombre omet réapparaît comme mémoire, incidence,
géométrie et structure physique.

% RESULTADO_RECUPERADO. Owner íntegro de la dinámica del cociente, el retorno
% observable y los cociclos de modo y orientación. Se subordina a la única
% construcción APP--TRIT--TPK sin abrir un origen alternativo.
\begin{HMTScientificOwnerSubordinated}
\input{sections/hmt/05b_extensiones_dinamica_cociente}
\end{HMTScientificOwnerSubordinated}
}
